\chapter{For loops}

\section{Brace expansion}

This topic should be its own chapter. Brace expansion happens when we use braces `\{\}' in bash.
You see braces are like wild cards and can help us create patterns to save ourselfs some effort. However,
in contrast to wild cards, they are not file and path dependent but will instead work with any string pattern.
Something to know is that brace expansions are simply to wide of a topic for us to check so we'll see some basic
uses of it and call it a day.

\lstinputlisting[
    language=bash,
    caption={11-bash.sh},
    captionpos=t,
    label={lst:11-bash},
    breaklines=true]
{../11-bash/11-bash.sh}

Something to consider is that bash expands the contents of the BE (because they work like preprocessors)
There is an order for which bash does it (since this process also happens with aliases, wild cards etc.).
For instance, variable expansion happens before brace expansion. This means that `\{1..\$int\}'
won't work. We might see a workaround for dynamic ranges later.

\section{Brace expansion in for loops}

Now lets get into what we are here for.


\lstinputlisting[
    language=bash,
    caption={11-bash.sh},
    captionpos=t,
    label={lst:11.1-bash},
    breaklines=true]
{../11-bash/11.1-bash.sh}

\section{Working with files}

This is bash, of course we can iterate through a list of files instead of a list of numbers.

\lstinputlisting[
    language=bash,
    caption={11-bash.sh},
    captionpos=t,
    label={lst:11.2-bash},
    breaklines=true]
{../11-bash/11.2-bash.sh}